\section{Causal mediation: proofs}
\label{sec:annex-p16-causal-mediation}

\begin{annexitem}{annex:p16.mediation-measurability}{thm:p16-mediation-measurability}{Composition before any probabilistic factorization}
\begin{proof}
In a finite model, first evaluate \(M_b(\omega)\), then evaluate the
already defined intervention table \(m\mapsto Y_{a,m}(\omega)\) at
that value. Both operations are total. In the \(A=a\) world, replacing
the mediator by the value it already attains preserves that world's
solution by \cref{thm:p16-counterfactual-laws}. Thus
\(Y_{a,M_a}(\omega)=Y_a(\omega)\). On the factual event \(A=a\),
the same theorem gives \(M_a=M\), and on \(A=a,M=m\) gives
\(Y_{a,m}=Y\). No division by the mass of either event was used.

For the general structural assertion, fix the intervened coordinate set
and let its assigned values vary in their finite product of standard
Borel spaces. Put the finite graph in topological order. The first
coordinate is either a projection onto a supplied intervention value or
a measurable function of the background. If all preceding coordinates
are measurable functions of the intervention values and background, the
tuple of its parent coordinates is measurable for the product sigma-algebra.
Composing that tuple and its disturbance coordinate with the stipulated
jointly measurable structural function gives a measurable next coordinate.
Induction proves joint measurability of the whole unique solution; the
same topological induction proves uniqueness without a finiteness
assumption on the coordinate value sets. Jointly measurable substitution
of the random \(M_b\) into the pair-intervention solution now proves
measurability of the nested response. The same equation-comparison
argument proves consistency and composition pointwise. In the abstract
formulation they are explicit assumptions, not inferred from
measurability alone.
\end{proof}
\end{annexitem}

\begin{annexitem}{annex:p16.mediation-decompositions}{thm:p16-mediation-decompositions}{Telescoping the four nested responses}
\begin{proof}
Write \(Z_{ij}=Y_{a_i,M_{a_j}}\). Composition identifies \(Z_{ii}\)
with \(Y_{a_i}\). At each outcome,
\[
 Z_{11}-Z_{00}=(Z_{10}-Z_{00})+(Z_{11}-Z_{10})
             =(Z_{11}-Z_{01})+(Z_{01}-Z_{00}).
\]
These are the two stated decompositions. Subtracting either pure
component from the corresponding total component gives
\(Z_{11}-Z_{10}-Z_{01}+Z_{00}\). Integrability permits finite linearity
of expectation and gives all the mean identities. Reverse-transition
NDE is \(\theta_{01}-\theta_{11}=-\operatorname{TNDE}\), while
reverse-transition NIE is \(\theta_{10}-\theta_{11}=-\operatorname{TNIE}\).
Substitution gives exactly the two source identities.

Under the additive response restriction, the contrast in exposure at
either fixed mediator value is \(u(a_1,C,R)-u(a_0,C,R)\), so
the difference between that contrast at \(M_{a_1}\) and at \(M_{a_0}\)
is zero pointwise. This proves \(I=0\). Without this restriction,
\(I\) is the expectation of precisely that difference. The nonlinear
example below has nonzero \(I\), so the extra equality of the two
forward pure components with TE cannot be inferred by algebra alone.
\end{proof}
\end{annexitem}

\begin{annexitem}{annex:p16.mediation-countable-formula}{thm:p16-mediation-countable-formula}{The original summation route and its support}
\begin{proof}
For a bounded measurable \(\varphi\), pointwise finite or countable
partition by \(M_b\) gives
\[
 \varphi(Y_{a,M_b})=\sum_m\mathbf1_{\{M_b=m\}}\varphi(Y_{a,m}).
\]
For nonnegative \(\varphi\), monotone convergence permits both expectation
and conditional expectation through this sum. Apply the assumed
conditional independence at each fixed \(m\). The corresponding
conditional product expectation equals
\[
 \eta_b(C,\{m\})\int\varphi(y)Q_a(m,C,dy)
 \quad\hbox{almost surely}.
\]
There are only countably many mediator values, so the exceptional sets
can be joined into one baseline-null set. Integrating over the baseline
law gives (MF). Subtract positive and negative parts to handle bounded
signed functions. The same nonnegative calculation for absolute values
shows that the stated absolute integral is exactly the absolute
expectation of the nested response, and permits the mean formula under
the specified integrability condition.

Conditional probability kernels for each real \(Y_{a,m}\) exist by
\cref{thm:p15-value-conditional-law}. Countability of the mediator makes
their assembly a jointly measurable family. A countable generating class
for the response Borel sets and the countable mediator index also provide
one baseline-null set for simultaneous equality of all resulting measures.
If \(\eta_b(c,\{m\})=0\), this cell contributes zero before a conditional
response mean is queried. Knowledge of the two families of factors on
the remaining support therefore determines every bounded response
functional, and the mean when integrable.

Apply this result first at \((a_1,a_0)\), then at \((a_0,a_0)\), and
subtract to obtain the NDE expression. Apply it at \((a_0,a_1)\) and
\((a_0,a_0)\) to obtain PNIE. This explicitly states the diagonal
factorization needed to replace the baseline mean by its mediator
mixture; alternatively that single-world mean may be identified directly
from a separate exposure intervention. The factor-identification route
is independent of how those factors were obtained. Equality of identified
factors in two observationally compatible models, followed by this same
integral, proves Pearl's nonexperimental reduction as well as his
experimental reduction. The causal cross-world premise remains an assumption
of this route even when all ordinary experiments have supplied the factors.
\end{proof}
\end{annexitem}

\begin{annexitem}{annex:p16.mediation-kernel-formula}{thm:p16-mediation-kernel-formula}{One joint kernel licenses the random evaluation}
\begin{proof}
The regular conditional laws \(\kappa\) and \(\eta_b\) exist by the
standard Borel conditional-law theorem. For every Borel response set
\(B\), the integrand \((a,m,c,r)\mapsto\mathbf1_B(f(a,m,c,r))\)
is jointly measurable. The kernel section-integral theorem
\cref{thm:p15-kernel-joint-measure}, with \((a,m,c)\) as parameter and
kernel \(\kappa(c,dr)\), makes \(Q_a(m,c,B)\) jointly measurable.
At each parameter it is the pushforward of a probability measure under
a measurable map, hence a probability measure in \(B\).

By (CW), the law of \((C,M_b,R)\) agrees on every measurable rectangle
with \(\nu(dc)\eta_b(c,dm)\kappa(c,dr)\). The kernel joint-measure
construction gives this product a probability measure. Uniqueness on
the generating rectangle pi-system identifies it with the actual joint
law. Applying the kernel integration formula to the jointly measurable
function \(\varphi(f(a,m,c,r))\) yields
\[
\begin{aligned}
 \mathbb E[\varphi(Y_{a,M_b})]
 &=\int\!\int\!\int
       \varphi(f(a,m,c,r))\kappa(c,dr)\eta_b(c,dm)\nu(dc)\\
 &=\int\!\int\!\int\varphi(y)Q_a(m,c,dy)\eta_b(c,dm)\nu(dc).
\end{aligned}
\]
Nonnegative functions use monotone convergence; signed functions use
the stated absolute-integrability condition before subtraction.
The same calculation with an indicator of a baseline set proves that
\(c\mapsto\int Q_a(m,c,\cdot)\eta_b(c,dm)\) is a conditional law
of the nested response. It is a probability kernel by kernel integration.
The uniqueness theorem for real conditional laws gives a single
baseline-null set on which equality might fail as measures, by first
comparing rational half-lines and then their generated sigma-algebra.

This proof integrates the actual \emph{joint} law of the random mediator
and response background. It does not evaluate a collection of scalar
conditional versions along a diagonal. For a countable mediator, the
previous proof needed only individual-response conditional independence;
(CW) is a sufficient stronger condition here. Its failure to follow
from pointwise uncountable assertions is proved by the diagonal
counterexample below.
\end{proof}
\end{annexitem}

\begin{annexitem}{annex:p16.mediation-finite-observed}{thm:p16-mediation-finite-observed}{Every observational replacement has a positive denominator}
\begin{proof}
Fix a baseline cell \(c\) of positive mass and a target mediator value
with \(p(M=m\mid A=b,c)>0\). Both exposure options and the needed
mediator option under \(a\) are positive by hypothesis. Exposure
exchangeability and consistency give
\[
 \mathbb P(M_b=m\mid C=c)
 =\mathbb P(M_b=m\mid A=b,C=c)
 =p(M=m\mid A=b,C=c).
\]
For any bounded response test function, the other two ordinary
exchangeability assumptions and consistency give in succession
\[
\begin{aligned}
 \mathbb E[\varphi(Y_{a,m})\mid c]
 &=\mathbb E[\varphi(Y_{a,m})\mid A=a,c]\\
 &=\mathbb E[\varphi(Y_{a,m})\mid A=a,M=m,c]\\
 &=\mathbb E[\varphi(Y)\mid A=a,M=m,c].
\end{aligned}
\]
The first division uses positive exposure probability, the second
positive mediator probability, and the last equality is pointwise
consistency on the resulting positive event. Insert these factors into
the proved countable cross-world formula. For \(\varphi(y)=y\), this
gives (OMF). Zero baseline or reference-mediator cells contribute zero
and need no ratio. Omitting the cross-world step leaves an identification
of the controlled-response factors and the mediator marginal; the
independent-draw construction below then identifies their product
mixture without determining the natural coupling. For a fixed controlled
mean the same conditional-response calculation is integrated over \(c\),
so positivity must cover that one mediator value at every baseline stratum.

To exhibit failure if overlap is removed, let \(A\) be fair and
\(M=A\). Compare \(Y=M\) with \(Y=A\). Both observed laws put
mass \(1/2\) at \((A,M,Y)=(0,0,0),(1,1,1)\). Every fixed
potential response and mediator is constant, so all the independence
conditions hold. But \(Y_{1,M_0}\) is zero in the first model and one
in the second, giving different natural direct effects. The missing
cell is \(A=1,M=0\). Positivity cannot be replaced by an arbitrary
conditional version there.
\end{proof}
\end{annexitem}

\begin{annexitem}{annex:p16.mediation-general-observed}{thm:p16-mediation-general-observed}{Common conditional measures and absolute-continuity transfer}
\begin{proof}
We give the proof first for a bounded measurable response test function
\(\varphi\); write \(\mu_a^\varphi\) for its integral against the observed
conditional response kernel. The mean follows by the stipulated absolute integration.
All conditional distributions below are regular kernels on standard
Borel spaces. Fix \(a\) with \(e_a(c)>0\) for \(\nu\)-almost every
\(c\). Its unconditional probability is positive. The baseline law
given \(A=a\) is \(\nu_a(dc)=e_a(c)\nu(dc)/P(A=a)\), equivalent
to \(\nu\), since the density is positive almost everywhere.

Let \(\tau_a(c,l,dr)\) be a conditional law of \(R\) given
\((A=a,C=c,L=l)\). Uniform mediator exchangeability implies that
the law of \((C,L,M,R)\) in that exposure group is
\[
 \nu_a(dc)\lambda_a(c,dl)\rho_a(c,l,dm)\tau_a(c,l,dr).       \tag{1}
\]
This is equality of joint measures, obtained by the conditional product
identity and the kernel joint-measure theorem. Consistency is the
jointly measurable identity \(Y=f(a,M,C,R)\) in this group. Integrating
\(\varphi(f(a,m,c,r))\) against (1) shows that
\[
 \widetilde\mu_a^\varphi(c,l,m)
   =\int\varphi(f(a,m,c,r))\tau_a(c,l,dr)
\]
is a jointly measurable conditional-response expectation. For indicator
test functions these integrals form one probability kernel; conditional-law
uniqueness gives equality with the observed response kernel outside one
\(\nu_a\lambda_a\rho_a\)-null set, simultaneously for all Borel response
events. It follows for every bounded \(\varphi\).

The target measure \(\nu\lambda_a k\) is absolutely continuous with
respect to \(\nu_a\lambda_a\rho_a\). Indeed, if a measurable set has
zero mass under the latter measure, equivalence of \(\nu_a\) and
\(\nu\) and nonnegative integration show that its mediator section
has \(\rho_a(c,l,\cdot)\)-mass zero for \(\nu\lambda_a\)-almost every
\((c,l)\). Condition (OV) gives zero \(k(c,\cdot)\)-mass for those
sections. Integrating proves the asserted absolute continuity. Thus
the observed and structural conditional kernels agree also under the
target integral. This is the step that prevents arbitrary null versions
from influencing the answer.

Uniform exposure exchangeability similarly gives, as an equality of
conditional probability measures for almost every \(c\),
\[
 \int\tau_a(c,l,dr)\lambda_a(c,dl)=\kappa(c,dr).
\]
Because \(k\) depends on \(c\) and not \(l\), kernel integration and
the last identity give
\[
\begin{aligned}
 \int\!\int\!\int\mu_a^\varphi(c,l,m)k(c,dm)\lambda_a(c,dl)\nu(dc)
 &=\int\!\int\!\int\varphi(f(a,m,c,r))\kappa(c,dr)k(c,dm)\nu(dc).
\end{aligned}
\]
The right side is the independent-draw response expectation by its
product construction. The same argument with absolute values supplies
the signed mean result whenever the target absolute integral is finite.
With (CW) and \(k=\eta_b\), the preceding general mediation theorem
identifies this expectation with the natural nested response.

Finally \(M_b\perp A\mid C\), positive \(e_b\), and consistency
identify its conditional law with the observed mediator law in group
\(b\), first on a countable generating class and then as measures outside
one baseline-null set. This proves the claimed reference-kernel
identification. A fixed controlled mediator is the kernel \(\delta_m\).
For non-atomic exposure laws the equivalence of \(\nu_a\) and \(\nu\)
cannot even be formed through a positive group event at a fixed \(a\);
no fixed-exposure identification assertion has been made for that case.
\end{proof}
\end{annexitem}

\begin{annexitem}{annex:p16.mediation-four-set-criterion}{thm:p16-mediation-four-set-criterion}{The cross-world graph bridge and the two distinct adjustment calculations}
\begin{proof}
\paragraph{The first separation supplies cross-world independence.}
We justify this step explicitly because a graph about one world cannot
simply be relabeled as two worlds. Add an independent exogenous root for
each disturbance and its edge to the corresponding endogenous vertex.
This changes none of the separations among the old vertices: each added
root has only one child and cannot lie internally on a simple path between
old vertices. Work with this expanded DAG and put
\(H=G_{\underline A\underline M}\). Restrict to the ancestors of
\(\{Y\}\cup M\cup W_0\), moralize, and remove \(W_0\).
\Cref{thm:p16-moral-separation-bridge} implies that the component containing \(Y\)
is disjoint from every component meeting \(M\). Call the former
\(D_Y\) and the union of the latter components \(D_M\).

To preserve the joint natural mediator tuple, restore each original
outgoing edge from a member of \(M\) whose child is an ancestor of
some member of \(M\) in the graph with exposure inputs cut. This
restoration joins only mediator-side components. Indeed, from such a
child follow a directed route to the first mediator encountered; before
that endpoint the route uses no outgoing mediator edge, and so already
exists in \(H\). It cannot cross \(W_0\), whose vertices are not
descendants of any mediator. Thus the child lies in \(D_M\). Every
other remaining nonbaseline parent of that child in the restored graph
lies there too, since its edge is already in \(H\), or it is itself a mediator. Consequently any
new moral edge also joins only \(D_M\) or a removed baseline vertex.
No connection to \(D_Y\) is created. The ancestral vertex set has
not increased: every original ancestor route to the mediator set can
be stopped at its first mediator, giving a route in \(H\). Denote
the restored DAG by \(H'\). Its separation of \(Y\) and \(M\)
given \(W_0\) therefore persists.

Make an auxiliary structural model with the graph \(H'\): replace each
mediator input along an edge not restored by its corresponding component
of the constant tuple \(m\); retain all restored mediator inputs. Replace
every input from \(A\) by \(a\) if its child lies in \(D_Y\), and
by \(b\) otherwise. Retain the original equations at \(A,M\), except
for these replaced parent inputs, and retain all disturbance laws.
Vertices outside the ancestral graph can be treated by the same rule;
they do not affect its targets. All replacement inputs are constants,
so the auxiliary graph is a subgraph of \(H'\) with independent
disturbances.

Every ancestor of \(W_0\) is unaffected by these changes: otherwise
there would be a directed path from \(A\) or \(M\) to \(W_0\),
contrary to the nondescendant condition. Thus its baseline values agree
with the original model pointwise. Along a directed ancestral route
to \(Y\) in \(H\), stop whenever a vertex of \(W_0\) is reached.
Every remaining vertex on the route is in \(D_Y\), because its directed
edges are also moral edges that remain after baseline removal. Therefore
the auxiliary equations relevant to \(Y\), up to these unchanged
baseline boundary values, have exactly the exposure inputs \(a\) and
mediator inputs \(m\). Induction in topological order proves that the
auxiliary response is \(Y_{a,m}\).

The same reasoning on routes to the mediator set gives exposure inputs
\(b\). All mediator-to-mediator ancestral routes have been restored,
including routes through vertices outside the mediator set. Their
nonbaseline vertices are in \(D_M\), and none is in \(D_Y\).
Topological induction on these routes gives the entire auxiliary tuple
exactly as \(M_b\). This proves both pointwise
identifications on one original outcome space. Applying
\cref{thm:p16-finite-global-markov} in the auxiliary graph gives
\(Y_{a,m}\perp M_b\mid W_0\), for arbitrary constants \(a,b,m\).
This is the source's cross-world graph license with its missing comparison
step supplied.

\paragraph{Controlled-response identification.}
For the second graphical condition, replace each mediator equation by its
component of the constant tuple \(m\), retain the natural equation of \(A\),
and replace its outgoing inputs by \(a\). Because every mediator is a descendant of \(A\),
the natural exposure is unchanged. The resulting response is \(Y_{a,m}\)
and the unchanged covariates are \((W_0,W_1)\). The second separation
and finite global Markov therefore give
\(Y_{a,m}\perp A\mid W_0,W_1\).
For the third, retain the natural mediator tuple, replacing its outgoing
inputs by \(m\) except on the mediator ancestral routes just described.
Apply the same restoration argument to the third separation with
\((A,W_0,W_1,W_2)\) as the conditioning set; its vertices are
nondescendants of every mediator. The resulting response is \(Y_m\), and the variables
\((A,W_0,W_1,W_2)\) are unchanged, since none is a descendant of \(M\).
The third separation gives their conditional independence of \(Y_m\)
and \(M\). On \(A=a\), composition identifies \(Y_m\) with
\(Y_{a,m}\). This proves the corresponding mediator exchangeability.

Now condition on a positive \(W_0=c\). Average the controlled response
over \(W_1\), insert \(A=a\) by its proved independence, partition
by \(W_2\), insert \(M=m\) by its proved independence, and use
consistency. Written without any omitted averaging, these steps are
\[
\begin{aligned}
 E[Y_{a,m}\mid c]
 &=\sum_{w_1}E[Y_{a,m}\mid c,w_1]p(w_1\mid c)\\
 &=\sum_{w_1,w_2}E[Y_{a,m}\mid a,c,w_1,w_2]
                           p(w_2\mid a,c,w_1)p(w_1\mid c)\\
 &=\sum_{w_1,w_2}E[Y\mid a,m,c,w_1,w_2]
                           p(w_2\mid a,c,w_1)p(w_1\mid c).
\end{aligned}
\]
The positivity conditions in the statement license every inserted
conditional event; a zero covariate weight is omitted before insertion.
Overlaps between the covariate sets simply restrict the sums to consistent
assignments, and cause no extra stochastic coordinates.

\paragraph{Reference-mediator identification.}
Retain natural \(A\) while feeding \(b\) to all its children. The
resulting mediator is \(M_b\), and \((W_0,W_3)\) is unchanged.
The fourth separation gives \(M_b\perp A\mid W_0,W_3\).
Condition on \(c\), partition by \(W_3\), insert the supported
\(A=b\), and use consistency to obtain
\(\eta_b(c,m)=\sum_{w_3}p(m\mid b,c,w_3)p(w_3\mid c)\).
Both factors are now identified, and the first part of the proof licenses
their cross-world product. Substitution in (MF) concludes the argument.
\end{proof}
\end{annexitem}

\begin{annexitem}{annex:p16.mediation-parent-scope}{thm:p16-mediation-parent-scope}{The exact full-parent restriction in the Markovian corollary}
\begin{proof}
Once all response parents have been set, its equation is
\(Y_{a,m}=f_Y(a,m,U_Y)\). Every parent of \(Y\) is a strict
nondescendant of \(Y\) in an acyclic graph. The entire mediator tuple
under an exposure intervention is consequently a function of
disturbances other than \(U_Y\); there is no directed path from
\(Y\) back to one of its parents. Independence of disturbances makes
\(U_Y\) independent of this whole tuple, so the cross-world
factorization holds, jointly over all its coordinates.

The observed parent tuple is also a function of disturbances other than
\(U_Y\). Thus the conditional law of the response at a positive parent
stratum is exactly the pushforward of the unchanged \(U_Y\) law under
\(f_Y(a,m,\cdot)\). This identifies the controlled kernel. When the
assumed adjustment set \(S\) identifies the mediator law,
\(\eta_b(m)=\sum_s p(m\mid b,s)p(s)\); substitute that kernel into
the countable formula and interchange the finite sums. This is the
displayed formula. In the unconfounded case its assumed mediator law
is \(p(m\mid b)\), giving the further specialization. Each
observational denominator is needed only where its factor contributes
positive target mass; the hypothesis expressly requires those cells.

The response-background argument uses only independence and measurable
functions of disjoint disturbance groups, so with standard Borel
coordinates it gives the common product kernel (CW) with \(R=U_Y\).
The general kernel mediation theorem then applies. It does not by itself
identify the mediator intervention law from a non-atomic observational
exposure slice; that factor needs its own valid identification argument.
Nor does the proof work by merely selecting one of several response
parents as the mediator: a remaining parent can carry exposure-induced
dependence into both the mediator and response. The explicit pair below
demonstrates this failure in independent-error models.
\end{proof}
\end{annexitem}

\begin{annexitem}{annex:p16.mediation-policy-identities}{thm:p16-mediation-policy-identities}{Independent kernel draws and their actual total contrast}
\begin{proof}
The map \(\omega\mapsto k(C(\omega),B)\) is measurable for every
mediator event. Hence it is a probability kernel from the original
outcome space, and \cref{thm:p15-kernel-joint-measure} constructs the
stated enlarged measure. Its marginal on the original space is unchanged.
The response is a composition of jointly measurable maps, so is measurable.
The original joint law of \((C,R)\) is \(\nu(dc)\kappa(c,dr)\).
For each test function, integration under the enlarged measure therefore
gives
\[
 \mathbb E[\varphi(Y_{a,G^k})]
 =\int\!\int\!\int\varphi(f(a,m,c,r))k(c,dm)\kappa(c,dr)\nu(dc).
\]
Interchanging the two kernel integrals by their joint product measure
gives the claimed \(Q_a\) mixture. Here the independent mediator draw
was imposed in the measure construction, not inferred about \(M_b\).

Both interventional decompositions are ordinary telescoping identities
among four numbers. Under (CW), the natural formula gives \(J_{ij}
=\theta_{ij}\) for every required pair, which proves equality of the
corresponding contrasts. Without it, equality of marginal mediator laws
does not imply equality of their couplings with the response background.
The induced-confounding example calculates a nonzero difference between
the two overall contrasts.

For the limited finite diagonal assertion with constant \(L\), the
ordinary exchangeability conditions identify the controlled-response
factors and mediator law, so
\[
 J_{ii}=\sum_c\sum_m\mathbb E[Y\mid a_i,m,c]p(m\mid a_i,c)p(c)
       =\sum_c\mathbb E[Y\mid a_i,c]p(c).
\]
If \(Y_{a_i}\perp A\mid C\) is also assumed, its consistency and
positive exposure probability identify the last expression with
\(\mathbb E[Y_{a_i}]\). This additional single-world response
exchangeability is explicit in the statement; no off-diagonal
cross-world implication is used.
\end{proof}
\end{annexitem}

\begin{annexitem}{annex:p16.mediation-nonlinear-example}{thm:p16-mediation-nonlinear-example}{All four nested means and the controlled contrast}
\begin{proof}
The independent treatment and mediator noise give the four stated
observed masses by multiplication. Under exposure zero the mediator
success probability is \(1/4\); under exposure one it is \(3/4\).
The response has no nonconstant disturbance and its controlled values
are \(Y_{0,m}=2m\) and \(Y_{1,m}=1+5m\). Thus
\[
 \theta_{00}=2(1/4)=1/2,\quad
 \theta_{01}=2(3/4)=3/2,\quad
 \theta_{10}=1+5(1/4)=9/4,\quad
 \theta_{11}=1+5(3/4)=19/4.
\]
Every required exposure and mediator cell has positive observed mass.
The controlled response is constant at every \((a,m)\), so all its
ordinary and cross-world independence conditions hold. Exposure is
independent of the complete mediator potential pair. Hence inserting
these same controlled values and mediator probabilities into (OMF)
gives the direct structural calculation in every cell.

Subtraction gives TE \(19/4-1/2=17/4\), PNDE \(9/4-1/2=7/4\),
PNIE \(3/2-1/2=1\), TNDE \(19/4-3/2=13/4\), and TNIE
\(19/4-9/4=5/2\). Their interaction difference is
\(19/4-9/4-3/2+1/2=3/2\). The controlled contrast is
\((1+5m)-2m=1+3m\). Its average over the reference mediator law
is \((3/4)1+(1/4)4=7/4\), the natural direct effect.
As a particularly simple limiting interaction example, the deterministic
unit \(M=A,Y=AM\) has TE one, both forward pure natural effects zero,
and controlled direct effect one at \(m=1\). These are structural
calculations; that limiting example lacks the mediator overlap of the
main four-cell example and is not offered as observational identification.
\end{proof}
\end{annexitem}

\begin{annexitem}{annex:p16.mediation-crossworld-countermodels}{thm:p16-mediation-crossworld-countermodels}{Positive finite laws and different natural couplings}
\begin{proof}
Let \(A,U,V\) be independent fair bits and \(E_0,E_1\) independent
Bernoulli\((1/4)\) bits independent of them. Set \(M_0=U,M_1=V\)
and \(M=(1-A)U+AV\). In model \(\mathcal M_+\) use
\[
 Y_{0,m}=E_0,\qquad Y_{1,m}=\mathbf1_{\{m=U\}}\oplus E_1;
\]
in model \(\mathcal M_-\) replace the indicator by
\(\mathbf1_{\{m=1-U\}}\). In both models set \(Y=Y_{A,M}\).
These are total finite structural equations with a shared background
\(U\) for mediator and response, and randomized \(A\).

In either model the full observed table is
\[
\begin{array}{cc|cc}
 A&M&p(Y=0,A,M)&p(Y=1,A,M)\\\hline
 0&0&3/16&1/16\\0&1&3/16&1/16\\
 1&0&1/8&1/8\\1&1&1/8&1/8
\end{array}
\]
At exposure zero, the independent \(E_0\) supplies the first two
rows. At exposure one, \(M=V\) is independent of the fair \(U\);
the indicator is fair at either fixed \(V\), and adding independent
noise preserves that fair conditional law. This proves the last two
rows, positivity, and normalization.

At any fixed mediator intervention, the response under exposure zero
is Bernoulli\((1/4)\), and under exposure one is fair in both models.
With mediator unmodified but exposure fixed, the same two conditional
response laws hold independently of the mediator, which is fair in
either exposure world. These facts prove equality of the full laws for
interventions on \(A\), on \(M\), on both, or on neither. If the
response is also fixed, the remaining \((A,M)\) law is the same in
both models: each unfixed exposure is fair, and each unfixed mediator
is fair conditional on either exposure; fixed coordinates take their
specified constants. This covers every subset of observed coordinates.

All potential variables are independent of randomized \(A\). For
\(a=0\), the controlled response \(E_0\) is independent of observed
\(M\) within \(A=0\). For \(a=1\), its \((U,E_1)\)-based controlled
response is independent of \(M=V\) within \(A=1\). Thus the three
ordinary finite conditions hold, as does mediator overlap.
The actual nested-response mean tables are
\[
\begin{array}{c|cccc}
 &\theta_{00}&\theta_{01}&\theta_{10}&\theta_{11}\\\hline
 \mathcal M_+&1/4&1/4&3/4&1/2\\
 \mathcal M_-&1/4&1/4&1/4&1/2
\end{array}
\]
Indeed substituting \(M_0=U\) makes the first model's exposure-one
indicator one and the second model's zero; substituting \(M_1=V\)
makes either indicator fair. Therefore PNDE is \(1/2\) or zero,
whereas TE is \(1/4\) in both. The controlled mean at exposure one
is \(1/2\) for either mediator value; at exposure zero it is \(1/4\).
Their common observed mediation expression gives PNDE \(1/4\).
An independently drawn fair reference mediator gives that very contrast,
so the interventional interpretation is identified.
Cross-world independence fails explicitly: for fixed \(m\), the
exposure-one success probability conditional on \(M_0=m\) is \(3/4\)
in \(\mathcal M_+\) and \(1/4\) in \(\mathcal M_-\), rather than
the unconditional value \(1/2\). This isolates the missing assumption
while retaining every ordinary experiment and support condition.
\end{proof}
\end{annexitem}

\begin{annexitem}{annex:p16.mediation-induced-confounding}{thm:p16-mediation-induced-confounding}{A coupling of the intermediate confounder's two worlds}
\begin{proof}
Let the independent disturbance of \(L\) be a pair \((L_0,L_1)\)
with masses parameterized by \(t\):
\[
\begin{array}{cc|c}
 L_0&L_1&\text{mass}\\\hline
 0&0&1/4+t\\0&1&1/2-t\\1&0&1/4-t\\1&1&t
\end{array}
\qquad t\in\{0,1/4\}.
\]
Each is a probability law, with \(P(L_0=1)=1/4\) and
\(P(L_1=1)=1/2\). Take fair randomized \(A\) and independent
\(E_M,E_Y\), both Bernoulli\((1/4)\), independent of that pair.
Define \(L=L_A\), \(M=L\oplus E_M\), and \(Y=(LM)\oplus E_Y\).
The disturbance vectors of \(A,L,M,Y\) are mutually independent;
their finite acyclic graph is the one in the statement.

Put \(p_0=1/4,p_1=1/2\) and \(q(0)=3/4,q(1)=1/4\).
Every observed cell has the explicit positive mass
\[
 p(a,l,m,y)=\tfrac12 p_a^l(1-p_a)^{1-l}
                         q(m\oplus l)q(y\oplus lm).
\]
Summing successively over \(y,m,l,a\) gives one. This full sixteen-cell
law is independent of \(t\). Every observed-coordinate intervention
also has the same law: if \(L\) is fixed, its disturbance is unused;
otherwise at any exposure its relevant coordinate has the same marginal
\(p_a\) in both models. The remaining independent equations and
disturbance laws coincide, and fixing any further subset simply replaces
the corresponding factors by constants. Randomized unfixed exposure
mixes those same exposure-specific laws. This verifies all subsets,
not only the observational law.

For any binary \(B\), independent \(E_Y\) gives
\(E[B\oplus E_Y]=1/4+(1/2)P(B=1)\). Also
\[
 E[L_i(L_j\oplus E_M)]=\tfrac14P(L_i=1)
                              +\tfrac12P(L_i=1,L_j=1).
\]
Consequently
\[
 \theta_{00}=11/32,\quad \theta_{11}=7/16,
 \quad\theta_{10}=5/16+t/4,\quad\theta_{01}=9/32+t/4.
\]
At \(t=0\) and \(t=1/4\), PNDE is respectively \(-1/32\)
and \(1/32\), while TE is \(7/16-11/32=3/32\).
PNIE is respectively \(-1/16\) and zero; TNIE is respectively
\(1/8\) and \(1/16\); TNDE is respectively \(5/32\) and \(3/32\).
Both decompositions can now be checked directly in either model.
The controlled direct effect is \(m/8\).

The two exposure-specific mediator success probabilities are
\(1/4+(1/2)p_0=3/8\) and \(1/4+(1/2)p_1=1/2\).
For \(m=0\) the observed conditional response mean is \(1/4\)
in either exposure group. For \(m=1\), Bayes' formula gives
\(P(L=1\mid A=0,M=1)=1/2\) and
\(P(L=1\mid A=1,M=1)=3/4\). Their conditional response means
are \(1/2\) and \(5/8\). The unadjusted mediator formula therefore
returns \((5/8-1/2)(3/8)=3/64\), different from either natural PNDE.

To compute the independent-draw target correctly, exposure \(a_i\)
leaves \(L_i\) with success probability \(p_i\), while the new mediator
draw has success probability \(r_j=1/4+p_j/2\), independently of
\(L_i\). Hence
\(J_{ij}=1/4+(1/2)p_i r_j\), giving exactly the displayed table.
Equivalently the observed conditional mean given \((a,l,m)\) is
\(1/4+lm/2\). Summing it against \(p(l\mid a_i)\) and the independent
reference mediator law gives the same expression (IGF). Every mediator
option given \((a,l)\) has probability \(1/4\) or \(3/4\), so the
required overlap holds. A response background \(R=(L_0,L_1,E_Y)\)
is independent of exposure and, conditional on \((A,L)\), independent
of observed \(M\), whose remaining randomness is \(E_M\). Thus the
uniform observational hypotheses are satisfied.

The interventional direct and indirect contrasts are \(3/64\) and
\(2/64\), so their sum is \(5/64\), different from the natural
total effect \(6/64\). If the selected mediator is enlarged to the
pair \((L,M)\), both other parents of \(Y\) are fixed at their
reference values and the response equation ignores \(A\). The natural
direct effect for that new mediator tuple is zero pointwise. This
proves both the original source's scope qualification and the
nonidentification of the smaller-mediator natural target.
\end{proof}
\end{annexitem}

\begin{annexitem}{annex:p16.mediation-diagonal}{thm:p16-mediation-diagonal}{A null set for every fixed mediator and a nonnull diagonal}
\begin{proof}
Let \(U,V\) be independent uniform variables on \([0,1]\), and
let \(A\) be an independent fair bit. In the first model set
\[
 M_0=U,\quad M_1=V,\quad M=(1-A)U+AV,
 \qquad Y_{0,m}=0,\quad Y_{1,m}=\mathbf1_{\{m=U\}},
 \quad Y=Y_{A,M}.
\]
In the second model retain the same mediator equations but make every
response zero. All spaces are standard Borel, and the diagonal
\(\{(m,u):m=u\}\) is closed in \([0,1]^2\), so these are jointly
measurable total structural equations. They satisfy composition and
consistency pointwise.

For every fixed \(m\), \(P(U=m)=0\). Thus every controlled response
in the first model is zero almost surely, exactly as in the second.
Such a response is independent of every random variable, including
\(A\), observed \(M\) in either exposure stratum, and every \(M_b\).
Exposure is independent of both mediator potentials. Hence all four
fixed-mediator independence conditions hold. Both observed mediator
laws given exposure are uniform, so the target/reference mediator laws
are mutually absolutely continuous.

At exposure one with natural mediator, the first response is
\(\mathbf1_{\{U=V\}}\). To see that its probability is zero, partition
\([0,1]\) into \(n\) equal half-open intervals, assigning the endpoint
one to the last interval. Apart from probability-zero interval boundaries,
\(U=V\) requires both variables to lie in the same interval, whose
total probability is \(n(1/n)^2=1/n\). Let \(n\) increase. Thus both
observed laws have fair independent exposure and uniform mediator with
response zero almost surely; the same holds under either exposure
intervention. Every controlled exposure/mediator intervention also has
response zero almost surely and hence agrees between the models.

But the first model has \(Y_{1,M_0}=\mathbf1_{\{U=U\}}=1\)
pointwise, while \(Y_0=0\). Its PNDE is one. The second has PNDE zero.
The separate fixed-mediator controlled kernels can all be chosen to be
\(\delta_0\), giving a mediator mixture of zero in either model.
No choice of finitely or countably many exceptional baseline sets repairs
that false conclusion: the offending null set varies with the
uncountable mediator index. In the first model the response background
\(R=(U,V)\) is not independent of \(M_0=U\); their joint law is
concentrated on equality of that coordinate. The required product law
(CW) therefore fails. This proves the exact need for a uniform
joint-kernel premise, or another explicitly justified regularity condition,
in the general-space replacement theorem.
\end{proof}
\end{annexitem}

\begin{annexitem}{annex:p16.mediation-path-semantics}{thm:p16-mediation-path-semantics}{Reference inputs retain the same unit's background}
\begin{proof}
First solve the original model with reference exposure \(a_0\).
Every reference coordinate is a measurable function of the original
background by the structural measurability theorem. At every modified
equation the retained current inputs belong to a subset of the original
parents, and all other inputs are these already computed background
functions. Thus the original topological order still evaluates the
modified equations uniquely. The same measurable-composition induction
proves joint measurability. Dependence introduced among the effective
background inputs is retained, even if the original disturbances were
independent.

Under reference exposure, induction shows that each current parent
input equals its reference value. Every retained or replaced input
therefore equals the input in the original reference solution, and the
response is \(Y_{a_0}\). In the source's Figure~3 the retained edges
give successively \(M=f_M(a_1,U_M)\),
\(L=f_L(M,a_0,U_L)\), and \(Y=f_Y(M_{a_0},L,U_Y)\).
Subtracting the reference response is exactly the formula in the body.

For its nonidentification claim use the preceding two finite models,
identifying the source's first intermediate \(M\) with their \(L\)
and its second intermediate \(L\) with their \(M\). Permit an additional
\(A\to M\) edge in their graph, ignored by the mediator equation;
the source allows an unused parent. The selected path sends \(L_1\)
through the mediator equation while the response's direct \(L\) input
is held at \(L_0\). The modified response is
\((L_0(L_1\oplus E_M))\oplus E_Y\), with mean
\(9/32+t/4\). Its reference mean is \(11/32\). The path-specific
effect is therefore \(-1/16\) when \(t=0\), and zero when \(t=1/4\).
All observed and single-world intervention laws coincide by the earlier
proof. The source's path-specific target consequently is not determined
by those laws even in the indicated independent-error class.
\end{proof}
\end{annexitem}
